\section{Tones, the Rule, and Emergent Time}
\label{sec:dynamics}

On the crystal sit the tones, moved by one local rule. This section covers the dynamics and the
emergent notions of energy and time (P1 to P6, P37, P44).

\paragraph{Tones and the rule.}
Each vibe carries a ternary tone $\{-1, 0, +1\}$, read as pain, peace, pleasure. Each note carries
a ternary fill, the sign and strength of attention. The update is signed majority: a vibe's next
tone is the sign of the sum, over its notes, of the neighbour's tone times the fill. There are no
continuous weights and no hidden state, only the three values and a transient integer tally. The
updates are local and asynchronous, with no global clock, since a universal now would be a preferred
frame.

\paragraph{The law as a trilemma, resolved (P3, P4).}
A local update on a graph faces a trilemma between regularity, Lorentz safety, and a sensible
energy. The resolution is the emergent-mesh Hamiltonian: on the random hyperbolic substrate the
update has an energy functional that is local, bounded below so the world does not fall forever, and
finite-speed, giving a light cone. P3 and P4 establish this Hamiltonian and its boundedness, and the
testbed reads its spectrum off the mesh Laplacian.

\paragraph{Emergent geometry and time (P5, P38).}
Space is the large-scale shape of the relations. P5 recovers a definite emergent dimension from the
graph alone, and P38 sharpens the emergent spatial geometry, with the recovered dimension stabilising
as the mesh grows (made unbiased in P69). Time is the depth of the before-and-after order, the length
of the longest causal chain. The arrow of time is the one-way accumulation of relations: the mesh
only grows, the interior is never edited, and the present is its leading edge.

\paragraph{The signal sector (P37).}
P37 shows that signals, propagating disturbances with a finite speed, come from the rule itself on
the grown mesh, not from any added wave equation. A local perturbation spreads at a bounded rate set
by the light cone, the discrete analogue of a field excitation.

\paragraph{Computational universality (P44).}
The single rule is computationally universal. With suitable fills it realises the basic logical
operation from which all circuits are built, and combined with the crystal's unlimited addressable
space it can host any pattern that can exist, including a working mind. This is the floor under the
claim that the substrate can express everything, kept separate from the harder question of felt
experience.
